
The spurious correlation problem in supervised learning — where models exploit features that correlate with labels in training data but not under distribution shift — has received sustained attention over the past several years. Benchmark datasets such as Waterbirds \citep{Sagawaetal2020} and CelebA \citep{Liuetal2021} have enabled systematic evaluation of mitigation strategies including Group Distributionally Robust Optimization (GroupDRO) \citep{Sagawaetal2020}, Just Train Twice (JTT) \citep{Liuetal2021}, and Deep Feature Reweighting (DFR) \citep{Kirichenkoetal2022}. These methods improve worst-group accuracy (WGA), in some cases substantially, without requiring group annotations at training time. DFR in particular proceeds by retraining only the final linear layer of an ERM-trained backbone using a class-balanced held-out split, relying on the observation that ERM backbones already encode core features \citep{Kirichenkoetal2022}.

Despite this empirical progress, the temporal dynamics underlying spurious feature acquisition remain unquantified. Existing methods exploit the consequences of spurious feature learning without measuring the process itself. Mangalam and Girshick [2021] observe qualitatively that shortcut features emerge preferentially during early training, but no systematic protocol has been established for measuring when this occurs, how large the temporal gap is, or how consistently it appears across random seeds. Prior theoretical work on the Frequency Principle \citep{Xuetal2019} and Simplicity Bias \citep{Shahetal2020} provides a mechanistic basis for expecting lower-complexity features to be learned first, but neither work provides an operationalization for standard spurious correlation benchmarks.

The present work addresses this gap by introducing the δ(t) framework: a protocol based on checkpoint linear probing that makes the temporal competition between spurious and core feature encoding directly measurable. At each training checkpoint, lightweight linear classifiers (probes) are trained separately for spurious and core labels on frozen backbone representations. Their accuracy difference, δ(t), constitutes a time series over training that captures which feature type dominates the backbone's representations at each epoch. The first epoch at which this difference becomes and remains negligible is designated t\*.

This paper reports a proof-of-concept application of the δ(t) framework to Waterbirds/ResNet-50 with three random seeds and 30 training epochs. The study is organized around five sub-hypotheses (H-E1, H-M1, H-M2, H-M3, H-M4) forming a sequential causal chain from feature complexity to gradient dynamics to temporal gap to DFR efficacy. Results are reported for each sub-hypothesis in turn, with explicit gate evaluations and limitation statements.

Contributions of this work are:

\begin{enumerate}
\item \textbf{Measurement framework:} A reproducible protocol for computing δ(t), gap area A, and transition epoch t\* on standard spurious correlation benchmarks, implemented and validated on Waterbirds/ResNet-50.
\end{enumerate}

\begin{enumerate}
\item \textbf{Empirical causal chain (partial):} Quantitative confirmation of the complexity hierarchy (H-M2: 3/3 metrics, all p<0.05 uncorrected) and gradient asymmetry (H-M1: GDR=6.977, Wilcoxon underpowered) driving the temporal gap (H-E1: p=0.022), with t\* exhibiting low cross-seed variance (H-M3: std=2.0 epochs).
\end{enumerate}

\begin{enumerate}
\item \textbf{Negative finding on DFR mechanism (H-M4):} The hypothesized positive correlation between DFR improvement and post-t\\textit{ training depth is not observed (r=−0.815). DFR WGA is robustly high at all training depths including severely undertrained backbones (epoch 1: WGA=0.806), which is inconsistent with t\} functioning as a necessary threshold for DFR applicability.
\end{enumerate}

